% Derived skeleton; official class/style files are unmodified.
\documentclass[10pt]{article}
\usepackage{tmlr}
\usepackage{graphicx,hyperref,url,algorithm,algorithmic}
\usepackage{amsmath,amssymb}
\newtheorem{PWTheorem}{Theorem}
\newcommand{\AuthorPending}[1]{\textbf{AUTHOR INPUT PENDING:} #1}
\title{AUTHOR INPUT PENDING: Manuscript Title}
\author{\name AUTHOR INFORMATION PENDING\email author-input-pending\\\addr AUTHOR INFORMATION PENDING}
\def\month{MM}\def\year{YYYY}\def\openreview{AUTHOR INPUT PENDING: OpenReview link}
\begin{document}\maketitle
\begin{abstract}
AUTHOR INPUT PENDING: write a single-paragraph abstract from the actual research problem,
method or theoretical contribution, key evidence, and scope. This template contains
no research findings and cannot be submitted as a manuscript.
\end{abstract}
\section{Introduction}
\label{sec:introduction}
\AuthorPending{Research problem, closest work, exact gap, contributions and evidence.}
This project is a formatting skeleton, not a research result. The citation
\cite{formatguide} points to the official author guidance solely to demonstrate bibliography formatting.

\section{Method and Analysis}
\label{sec:method}
\AuthorPending{Definitions, design decisions, assumptions, algorithm inputs and outputs.}
The following identity is a typesetting demonstration only:
\begin{equation}
\label{eq:identity}
\begin{aligned}
 (a+b)^2 &= a^2 + 2ab + b^2,\\
 \sum_{i=1}^{n}(x_i-\bar x)^2
 &= \sum_{i=1}^{n}x_i^2-n\bar x^2.
\end{aligned}
\end{equation}
\begin{PWTheorem}[Typesetting demonstration]
For a real number $x$, the square $x^2$ is nonnegative.
\end{PWTheorem}
\textit{Proof (format example).} The statement follows from the ordering of the real numbers.
The displayed theorem is not a contribution of the manuscript.

\begin{algorithm}[htbp]
\caption{Pseudocode formatting example only}
\label{alg:format}
\let\AND\relax
\begin{algorithmic}[1]
\REQUIRE Author-supplied inputs
\ENSURE Author-supplied outputs
\STATE Replace this line with the actual method.
\end{algorithmic}
\end{algorithm}

\section{Experiments or Further Analysis}
\label{sec:results}
\AuthorPending{Provide real evidence and conditions. Replace this section for theory-only papers.}
Equation~\ref{eq:identity}, Figure~\ref{fig:single}, Figure~\ref{fig:wide}, and
Table~\ref{tab:conditions} demonstrate cross-references, not measurements.
\begin{figure}[htbp]
\centering
\fbox{\parbox{0.86\linewidth}{\centering DRAFT PLACEHOLDER: insert a real figure with a source and meaningful caption.}}
\caption{Figure formatting example; no experimental result is represented.}
\label{fig:single}
\end{figure}
\begin{figure*}[t]
\centering
\fbox{\parbox{0.90\textwidth}{\centering DRAFT PLACEHOLDER: wide figure across the available page width.}}
\caption{Wide figure formatting example; replace with author-provided evidence.}
\label{fig:wide}
\end{figure*}
\begin{table}[htbp]
\centering
\caption{Reporting structure only; values have not been supplied.}
\label{tab:conditions}
\begin{tabular}{ll}
\hline
Item & Status\\
\hline
Data and task & Pending\\
Comparison conditions & Pending\\
Measurements and uncertainty & Pending\\
\hline
\end{tabular}
\end{table}

\section{Limitations and Conclusion}
\AuthorPending{State supported conclusions and claim-specific boundaries.}
\label{pw:main-end}

\bibliographystyle{tmlr}\bibliography{references}
\appendix
\section{Supporting Details}
\label{app:supporting}
\AuthorPending{Complete derivations, additional settings, and evidence referenced in the main text.}
This appendix is a formatting placeholder. It refers to Section~\ref{sec:method} and
Equation~\ref{eq:identity}; no omitted research result has been assumed.
\end{document}
